\section[Cote 154 : pseudo-droites, générisations et balayages]{Cote 154 : pseudo-droites, générisations et balayages\protect\verifie{Folder154}}\label{sec:154}

La cote 154, \q{[Système de pseudo-droites]} (notes manuscrites, 1983--1984), étudie les positions relatives d'une pseudo-droite ajoutée à un système $\Sigma$ de pseudo-droites d'un plan projectif réel : leurs spécialisations et générisations, la surface cellulaire qu'elles forment, la monodromie des polygones qu'elles découpent, et, dans le lot~8, des balayages de l'arrangement, écrits en diagrammes de fils, en tableaux de permutations et en mots de générateurs. Mathlib n'a ni les arrangements de pseudo-droites ni les matroïdes orientés : les énoncés vérifiés ici sont posés dans des modèles finis écrits pour l'occasion, ou, pour le dernier, dans le plan réel. Les trois premiers encadrés sont des pistes de la lecture ; les deux derniers sont tirés des dessins du lot~8, redessinés à la machine (\texttt{transcripts/154/drawings/}).

\begin{encadre}
\begin{conjecture*}[Grothendieck, page 18, l'ordre cyclique détermine la position\piste{154-cyclic-order-determines-position}]
Une position relative d'une pseudo-droite ajoutée, autre que les $D_i$ eux-mêmes, est déterminée par l'ordre cyclique, compatible à l'antipodie, dans lequel son revêtement double rencontre les pseudo-droites orientées de~$\Sigma$.
\end{conjecture*}
\end{encadre}
{\small La page écrit \q{il ne devrait pas être difficile de montrer} et ne le démontre pas. Mathlib n'a pas les arrangements de pseudo-droites. On se reportera aux notes de Grothendieck (\lot{154}{1}{18}).\par}

\begin{encadre}
\begin{theoreme*}[Grothendieck, pages 55 à 63, 85 et 86, monodromie diédrale\piste{154-dihedral-monodromy-half-turn}]\nonverifie
Soit $\Sigma$ un système de $n$ pseudo-droites d'un plan projectif réel. Les $2n$-gones $\mathrm{Pol}(D)$ découpés sur le revêtement double d'une pseudo-droite ajoutée $D$ forment un système local sur la surface cellulaire des positions relatives. Le transport le long d'un demi-tour autour d'une position $D = D_i$ est la rotation $t \mapsto t - (n-1)$ de $\mathbf Z/2n\mathbf Z$, et la monodromie $\pi_1 \to \mathbb D_{2n}$, à valeurs dans le groupe diédral d'ordre $4n$, est surjective sur la surface où la traversée des $D_i$ est permise.
\end{theoreme*}
\end{encadre}
{\small L'édition corrige deux fois $\mathbf Z/n\mathbf Z$ en $\mathbf Z/2n\mathbf Z$ (pages 57 et 85), et note que pour $n$ impair la surjectivité demande une réflexion latérale. La surjectivité sans traversée des $D_i$ (pages 85 et 86) n'est pas reprise. Mathlib n'a pas les arrangements de pseudo-droites ni les systèmes locaux sur un complexe cellulaire. On se reportera aux notes de Grothendieck (\lot{154}{3}{55}).\par}

\begin{encadre}
\begin{theoreme*}[Grothendieck, page 153, cube des générisations, sous la forme précisée par l'édition\piste{154-generisation-cube}]
Soit $\Sigma$ un arrangement de pseudo-droites, vu comme matroïde orienté de rang~$3$, et appelons \emph{position} d'une pseudo-droite ajoutée une signature $\sigma$ des cocircuits (les sommets de $\Sigma$, au signe près) valide sur chaque droite de $\Sigma$ au sens de Las Vergnas : le long de chaque droite, lue sur son revêtement double, les signes sont de la forme $+ \cdots + \,(0)\, - \cdots -$ au renversement près, avec au plus un $0$, et pas tous nuls. Disons que $\sigma$ est une générisation de $\sigma_D$ si $\sigma = \sigma_D$ partout où $\sigma_D \neq 0$. Si $\sigma_D$ s'annule en exactement $\nu$ sommets, les générisations de $D$ sont exactement les $3^\nu$ choix de signes, $0$ compris, en ces sommets ; ordonnées par la spécialisation, elles forment l'ensemble ordonné des faces du cube $[-1, 1]^\nu$. Les $2^\nu$ générisations stables (sans zéro) en sont les sommets, les $\nu\,2^{\nu-1}$ générisations sous-stables (un seul zéro) les arêtes, l'incidence étant la spécialisation.
\end{theoreme*}
\end{encadre}
{\small \emph{Démonstration.} Le point est que la condition de Las Vergnas se vérifie droite par droite, et qu'une droite ne porte jamais deux zéros de $\sigma_D$ : deux sommets de $D$ ne sont jamais sur une même droite de $\Sigma$. Sur une droite qui porte un zéro de $\sigma_D$, la suite des signes est $a \cdots a\, 0\, (-a) \cdots (-a)$ ; remplacer ce $0$ par $a$, par $-a$ ou le garder donne encore une suite valide. Sur une droite qui n'en porte pas, rien ne change. Tout choix de signes aux $\nu$ zéros de $\sigma_D$ est donc une position, et les choix aux différents sommets sont indépendants : c'est le \q{décollement} de la page, de chaque côté, en chaque sommet (figure~\ref{fig:154-cube}). La restriction aux zéros de $\sigma_D$ est alors une bijection des générisations de $D$ sur $\{-, 0, +\}^\nu$, qui respecte la spécialisation ($0$ en dessous de $+$ et de $-$). Le vecteur de signes $w$ correspond à la face du cube où la coordonnée $z$ vaut $w_z$ si $w_z \neq 0$ et parcourt $[-1, 1]$ sinon ; une face est contenue dans une autre exactement quand le vecteur de la seconde est une spécialisation de celui de la première, et les vecteurs sans zéro donnent les sommets. Il reste à compter : $2^\nu$ vecteurs sans zéro, et, pour chacun des $\nu$ emplacements du zéro unique, $2^{\nu-1}$ choix de signes ailleurs. L'hypothèse $\nu \geq 2$ n'est pas utile, ni les axiomes d'un arrangement (deux droites se coupent en un sommet, \dots) : l'énoncé vaut pour toute donnée de droites portant des suites de sommets. Lean (\href{\leanblob/Folder154.lean}{\texttt{Folder154.\allowbreak{}lean}}) définit les arrangements (\texttt{Arrangement}), la condition sur une droite (\texttt{ValidSeq}), les positions (\texttt{IsPosition}) et la spécialisation (\texttt{Spec}), et démontre qu'une droite porte au plus un zéro (\texttt{IsPosition.\allowbreak{}eq\_of\_zero\_on\_line}), l'indépendance (\texttt{IsPosition.\allowbreak{}of\_spec}, \texttt{isPosition\_and\_spec\_iff}), la bijection et l'incidence (\texttt{genEquiv}, \texttt{genEquiv\_spec\_iff}), le cube (\texttt{cubeFace\_subset\_iff}, \texttt{cubeFace\_eq\_singleton}) et les comptes (\texttt{card\_gen}, \texttt{card\_gen\_stable}, \texttt{card\_gen\_substable}). La page énonce le résultat sans justification ; l'indépendance est celle que proposait la lecture. Ne sont pas démontrés : le même énoncé pour des pseudo-droites topologiques, qui demanderait le théorème de représentation de Folkman--Lawrence ; la tension avec la page~150 pour $\nu \geq 3$ (les $2\nu$ positions sous-stables incidentes à un sommet de la surface $X^*$ contre les $\nu\,2^{\nu-1}$ arêtes du cube) ; le contraste avec une droite passant par $\nu$ sommets d'un arrangement de droites, qui n'a que $2\nu$ perturbations stables ; la question de la même page. On se reportera aux notes de Grothendieck (\lot{154}{8}{153}).\par}

\begin{figure}[htbp]
\centering
\begin{tikzpicture}[x=1cm, y=1cm, pt/.style={circle, fill, inner sep=1.2pt}, lab/.style={font=\scriptsize, inner sep=1pt}]
\coordinate (v1) at (1.2,0); \coordinate (v2) at (3.6,0);
\foreach \v in {v1,v2} { \draw[black!55] ($(\v)+(-0.8,-1.2)$) -- ($(\v)+(0.8,1.2)$); \draw[black!55] ($(\v)+(-0.8,1.2)$) -- ($(\v)+(0.8,-1.2)$); }
\draw[very thick] (-0.2,0) -- (5,0); \node[lab, right] at (5,0) {$D$};
\node[pt] at (v1) {}; \node[pt] at (v2) {};
\node[lab] at ($(v1)+(0,-0.45)$) {$v_1$}; \node[lab] at ($(v2)+(0,-0.45)$) {$v_2$};
\begin{scope}[xshift=8cm]
\fill[black!6] (-1.3,-1.3) rectangle (1.3,1.3);
\draw[thick] (-1.3,-1.3) rectangle (1.3,1.3);
\foreach \x/\sx in {-1/-,1/+} \foreach \y/\sy in {-1/-,1/+} { \node[pt] at ({1.3*\x},{1.3*\y}) {}; \node[lab] at ({1.75*\x},{1.6*\y}) {$(\sx,\sy)$}; }
\node[lab] at (0,1.55) {$(0,+)$}; \node[lab] at (0,-1.55) {$(0,-)$};
\node[lab] at (1.75,0) {$(+,0)$}; \node[lab] at (-1.75,0) {$(-,0)$};
\node[lab] at (0,0) {$(0,0)$};
\end{scope}
\end{tikzpicture}
\caption{Le cas $\nu = 2$. La pseudo-droite ajoutée $D$ passe par deux sommets $v_1$, $v_2$ de $\Sigma$. La décoller d'un côté ou de l'autre en chacun d'eux donne les quatre générisations stables, sommets du carré ; la décoller en un seul sommet donne ses quatre arêtes, et $D$ elle-même, de vecteur $(0, 0)$, est le carré entier (grisé). Figure de l'édition ; à la page~153 Grothendieck dessine un fuseau dont le sommet est traversé par $\Delta_1, \dots, \Delta_\nu$ (\lot{154}{8}{153}).}
\label{fig:154-cube}
\end{figure}

\begin{encadre}
\begin{theoreme*}[Grothendieck, pages 142, 144 et 145, un demi-balayage écrit deux fois\piste{154-sweep-142-144-same}]\verifie[\leanblobmain]{Folder154Sweep}
Numérotons les places de $1$ à $7$ et notons $\sigma_i$ l'échange des fils aux places $i$ et $i+1$. Le diagramme de fils de la page~142 fait croiser, de haut en bas, les paires $D/E$, $F/G$, $D/G$, $E/G$, $C/G$, $B/G$, $B/C$, $B/E$, $C/E$, $B/D$, $C/D$, $B/F$, $G/A$ : ce sont les croisements du mot $\sigma_4\sigma_6\sigma_5\sigma_4\sigma_3\sigma_2\sigma_3\sigma_4\sigma_3\sigma_5\sigma_4\sigma_6\sigma_1$, qui mène $A\,B\,C\,D\,E\,F\,G$ à $G\,A\,E\,D\,C\,F\,B$. Les dix lignes du tableau de la page~144 sont les états de ce balayage après $0$, $2$, $3$, $4$, $5$, $6$, $7$, $8$, $10$ et $13$ croisements : le diagramme et le tableau sont un même balayage. Le mot est réduit : ses $13$ lettres sont autant que les inversions de $G\,A\,E\,D\,C\,F\,B$, chaque paire croisée l'est une fois, et les huit paires jamais croisées sont $A/B$, $A/C$, $A/D$, $A/E$, $A/F$, $C/F$, $D/F$ et $E/F$. Le mot $\sigma_3\sigma_4\sigma_5\sigma_2\sigma_6\sigma_1\sigma_5\sigma_3$ des huit premiers termes de la page~145 est réduit ; avec le neuvième terme lu $\sigma_6$, les fils $C$ et $F$ se croisent une seconde fois et le mot cesse d'être réduit.
\end{theoreme*}
\end{encadre}
{\small \emph{Démonstration.} Un mot en les $\sigma_i$ a au moins autant de lettres que la permutation qu'il produit a d'inversions, car l'échange de deux places voisines ne change l'ordre que d'une paire de places (\texttt{inv\_mul\_adjSwap\_le}, \texttt{inv\_run\_le\_length}) ; un mot qui atteint cette borne est donc réduit, aucun mot plus court ne donnant la même permutation (\texttt{isReduced\_of\_inv\_eq}). Le reste est un calcul fini, que le noyau de Lean fait par \texttt{decide} : les paires croisées dans l'ordre (\texttt{pairs\_word142}), l'arrivée (\texttt{run\_word142}), les $13$ inversions (\texttt{inv\_run\_word142}, \texttt{isReduced\_word142}), les paires non croisées (\texttt{uncrossed142\_spec}), les lignes de la page~144 (\texttt{table144\_eq\_states}) ; pour la page~145, $8$ inversions pour $8$ lettres (\texttt{isReduced\_word145\_8}), la paire $C/F$ aux troisième et neuvième croisements (\texttt{word145\_9\_crosses\_CF\_twice}), et le mot $\sigma_3\sigma_2\sigma_1\sigma_4\sigma_3\sigma_6\sigma_5$, de sept lettres, qui donne la même permutation que les neuf premiers termes (\texttt{not\_isReduced\_word145\_9}). Entre deux lignes du tableau, les croisements du diagramme sont ceux que la ligne souligne, faits l'un après l'autre : $D/E$ et $F/G$ de $1'$ à $2'$, $C/E$ et $B/D$ de $8'$ à $9'$, $C/D$, $B/F$ et $G/A$ pour la dernière ligne ; ils portent sur des places disjointes et commutent. Les lignes du tableau et le mot de la page~145 sont ceux de la transcription, qui marque incertains le neuvième terme et les suivants, dans une tache d'encre (\lot{154}{8}{144}, \lot{154}{8}{145}). Les croisements de la page~142 sont lus sur le relevé 08-2, un redessin à la machine superposé au fac-similé, que personne n'a encore relu (\lot{154}{8}{142}). La lecture de la ligne $9'$ comme $A\,G\,E\,C\,D\,B\,F$ (la parenthèse $(A\,G)$ prise pour ces deux lettres dans cet ordre) et la traduction des croisements en mot sont de l'édition ; la convention qui fait agir $\sigma_i$ sur les places, et non sur les lettres, est de la lecture, et avec l'autre la vérification de la page~145 serait différente. L'énoncé de la fiche tient tel qu'il est écrit. N'est pas démontré : que les huit croisements du diagramme de fils des lignes $5'$ à $11'$ de la page~144 (relevé 08-18) soient $B/G$, $B/C$, $B/E$, $B/D$, $C/E$, $A/G$, $B/F$, $C/D$, ses fils n'étant pas identifiés à des lettres ; leur nombre seul est celui du tableau.\par}

\begin{encadre}
\begin{theoreme*}[D'après la page 153, l'ombre du cube est un $2\nu$-gone\pisteresolue{154-cube-shadow-2nu-gon}]\verifie[\leanblobmain]{Folder154Zonogon}
Soient $[p_1, q_1], \dots, [p_\nu, q_\nu]$, $\nu \geq 1$, des segments du plan $\mathbf R^2$, non réduits à un point, de directions deux à deux non parallèles. Leur somme de Minkowski $Z = [p_1, q_1] + \dots + [p_\nu, q_\nu]$ est symétrique par rapport au point $c = \sum_i (p_i + q_i)/2$ et a exactement $2\nu$ points extrémaux. En particulier, l'image du cube $[-1, 1]^\nu$ par une application linéaire vers le plan qui n'envoie pas deux vecteurs de base sur des vecteurs parallèles est un polygone convexe centralement symétrique à $2\nu$ sommets.
\end{theoreme*}
\end{encadre}
{\small \emph{Démonstration.} Posons $v_i = (q_i - p_i)/2$ : $Z$ est l'image du cube $[-1, 1]^\nu$ par $a \mapsto c + \sum_i a_i v_i$ (\texttt{sum\_segment\_eq}), et $a \mapsto -a$ donne la symétrie, sans hypothèse (\texttt{zono\_symm}, \texttt{sum\_segment\_symm}). Un point extrémal de $Z$ a toutes ses coordonnées égales à $\pm 1$ : si $|a_k| < 1$, il est le milieu des deux points $x \pm \delta v_k$ de $Z$ (\texttt{coeff\_of\_extreme}). Le sommet $x_\varepsilon = c + \sum_i \varepsilon_i v_i$ est extrémal si et seulement si une forme linéaire $f$ vérifie $\varepsilon_i f(v_i) > 0$ pour tout $i$ (\texttt{extremePoints\_zono}). Si $f$ existe, $x_\varepsilon$ est le seul point de $Z$ où $f$ atteint son maximum (\texttt{le\_ev\_vert}), donc il est extrémal, et deux vecteurs de signes réalisés distincts donnent deux points distincts (\texttt{vert\_injOn}). Sinon, le théorème de séparation (sous la forme de l'alternative de Gordan) fait de $0$ un barycentre $\sum_i \mu_i \varepsilon_i v_i$, à poids $\mu_i \geq 0$ de somme~$1$ ; alors $x_\varepsilon = c + \sum_i \varepsilon_i (1 - \mu_i) v_i$, et la coordonnée $\varepsilon_k (1 - \mu_k)$ d'un $k$ tel que $\mu_k > 0$ est strictement comprise entre $-1$ et $1$ : $x_\varepsilon$ n'est pas extrémal. Cette caractérisation vaut pour des $v_i$ non nuls quelconques ; seul le compte qui suit demande qu'ils ne soient pas parallèles. Les vecteurs de signes réalisés sont les régions de l'arrangement central des droites $f(v_i) = 0$ du plan dual, et il y en a $2\nu$ (\texttt{card\_R}). Par récurrence sur $\nu$ : pour $\nu = 1$, les deux signes sont réalisés par $\pm v_1$. Ajoutons un vecteur $w$. Toute région $t$ des autres vecteurs se prolonge, en déplaçant au besoin la forme le long de $w$ (\texttt{exists\_cons}, \texttt{both\_of\_zero}) ; elle se prolonge des deux façons si et seulement si elle contient $w^\perp$ ou $-w^\perp$, car le segment qui joint une forme réalisant $(t, +)$ à une forme réalisant $(t, -)$ reste dans la région $t$ et coupe la droite $f(w) = 0$ en un multiple non nul de $w^\perp$ (\texttt{eq\_of\_both}). Comme $w$ n'est parallèle à aucun $v_i$, les formes $\pm w^\perp$ ne s'annulent sur aucun $v_i$ et définissent deux régions distinctes (\texttt{pPlus\_ne\_pMinus}) : le compte augmente de~$2$ (\texttt{fibre\_card}). Le théorème réunit les deux parties (\texttt{ncard\_extremePoints\_sum\_segment}). La page écrit seulement, sous les comptes $2^\nu$ et $\nu\,2^{\nu-1}$ du cube des générisations, \q{Produit des diagonales d'un polygone convexe d'ordre $2\nu$}, et dessine au-dessous un contour fermé par huit points sans sommet intérieur, d'après le relevé 08-23, un redessin à la machine que personne n'a relu (\lot{154}{8}{153}). La lecture de cette ligne comme l'ombre du cube, l'énoncé et sa démonstration sont de l'édition. C'est un fait classique, les zonotopes étant les projections des cubes et leurs sommets répondant aux régions de l'arrangement central (G.~M.~Ziegler, \emph{Lectures on Polytopes}, leçon~7) : la piste est réglée, d'où le losange gris. Il est démontré ici pour tout $\nu$. Ne sont pas démontrés : l'identification des $2\nu$ sommets de l'ombre aux perturbations rectilignes stables d'une droite passant par $\nu$ sommets, les $2^\nu - 2\nu$ autres sommets du cube étant les générisations qu'aucune droite ne réalise, ni la lecture de la figure comme un octogone, c'est-à-dire $\nu = 4$.\par}
